\documentclass[11pt, a4paper]{article}

\usepackage[T1]{fontenc}
\usepackage[utf8]{inputenc}
\usepackage{tgpagella}
\usepackage{microtype}
\usepackage[spanish, es-noshorthands]{babel}
\addto\captionsspanish{\renewcommand{\tablename}{Tabla}}

\usepackage{amsmath, amssymb, bm}
\usepackage[table, xcdraw]{xcolor}
\usepackage{booktabs}
\usepackage{tabularx}
\usepackage[most]{tcolorbox}
\usepackage[margin=2.5cm]{geometry}
\usepackage{fancyhdr}

\pagestyle{fancy}
\fancyhf{}
\setlength{\headheight}{16pt}
\lhead{\textbf{UCB Sede Tarija} | Ing. Mecatrónica}
\chead{Solucionario y Diagnóstico (W08-S01)}
\rhead{Uso Docente}
\lfoot{Prof: M.Sc. Ing. Bernardo Quiroga Turdera}
\rfoot{Página \thepage}
\renewcommand{\headrulewidth}{0.4pt}
\renewcommand{\footrulewidth}{0.4pt}

\definecolor{ucbblue}{RGB}{0, 51, 102}

\begin{document}

\begin{center}
    \Large\textbf{\textcolor{red!80!black}{SOLUCIONARIO Y GUÍA DE DIAGNÓSTICO DOCENTE}}\\[0.2cm]
\end{center}

\section*{Solucionario de la Hoja de Trabajo[cite: 11]}

\textbf{Ejercicio 1:} $\dot{x} = -2\sin(q) \cdot \dot{q}$. El Jacobiano escalar es $J(q) = -2\sin(q)$.

\textbf{Ejercicio 2:}
$\mathbf{r}_i = \mathbf{o}_n - \mathbf{o}_{i-1} = [1,2,0]^T - [1,0,0]^T = [0,2,0]^T$.
Vel. Lineal: $J_{v_i} = [0,0,1]^T \times [0,2,0]^T = [-2,0,0]^T$.
Vel. Angular: $J_{\omega_i} = [0,0,1]^T$.
Total: $J_i = [-2, 0, 0, 0, 0, 1]^T$.

\textbf{Ejercicio 3:} Para una junta prismática $J_{\omega_i} = \mathbf{0}$. El eje aporta pura traslación, $J_i = [0, 1, 0, 0, 0, 0]^T$.

\textbf{Ejercicio 4:} Reemplazando $\sin(0)=0, \cos(0)=1, \sin(\pi/2)=1, \cos(\pi/2)=0$:
\begin{equation*}
    J = \begin{bmatrix} -1 & -1 \\ 1 & 0 \\ 0 & 0 \\ 0 & 0 \\ 0 & 0 \\ 1 & 1 \end{bmatrix}
\end{equation*}

\textbf{Ejercicio 5:}
$\mathbf{v}_e = J \begin{bmatrix} 0 \\ 2 \end{bmatrix} = \begin{bmatrix} -2 \\ 0 \\ 0 \\ 0 \\ 0 \\ 2 \end{bmatrix}$. 
Interpretación: Velocidad lineal de $-2$ en X, y el efector está rotando globalmente a $2$ rad/s en Z.

\textbf{Ejercicio 6:} No. El producto cruz requiere que ambos vectores operen en el mismo espacio vectorial (mismo marco de referencia). Operar mezclando marcos produce un vector basura sin sentido físico.

\textbf{Ejercicio 7:} Las raíces de $\sin(q_2) = 0$ en el rango dado son $q_2 = 0$ y $q_2 = \pi$.

\textbf{Ejercicio 8:} El primer mapeo es IK Global (Analítica), resuelve la pose. El segundo es IK Diferencial (Local), requiere conocer un $\mathbf{q}$ actual para invertir el Jacobiano y generar una velocidad continua $\dot{\mathbf{q}}$. El $J^{-1}$ no reemplaza a la IK Global (Pieper) porque no puede encontrar los ángulos iniciales partiendo de cero.

\vspace{0.5cm}

\section*{Guía de Intervención Diagnóstica para el Aula[cite: 11]}
\renewcommand{\arraystretch}{1.5}
\begin{tabularx}{\textwidth}{|X|X|X|}
\hline
\rowcolor{gray!20} \textbf{Observación / Error} & \textbf{Concepto Probablemente Débil} & \textbf{Intervención Inmediata} \\ \hline
Asume que el Jacobiano es una matriz de transformación homogénea. & Confusión de objetos matemáticos. & Contrastar matrices: $T$ dimensiona posiciones $4\times4$; $J$ proyecta velocidades $6\times n$. \\ \hline
Calcula el producto cruz en reversa $(\mathbf{r} \times \mathbf{z})$. & Dirección de velocidad tangencial errónea. & Recuerde la física de cuerpo rígido $\mathbf{v} = \boldsymbol{\omega} \times \mathbf{r}$. \\ \hline
Coloca componente angular no nula en columna prismática. & Naturaleza física de la junta no comprendida. & Contraste visualmente que un actuador lineal no hace rotar a su efector propio. \\ \hline
Evalúa $J\dot{q}$ como si fuera un punto en el espacio. & Comprensión de unidades y física ausente. & Requiera explícitamente escribir las unidades ($\mathrm{m/s}$ y $\mathrm{rad/s}$) al final. \\ \hline
Afirma que en singularidad el robot no se mueve en absoluto. & Naturaleza matricial de la pérdida de rango. & Muestre que las columnas del Jacobiano aún operan (el brazo "bare" se mueve en Y tangencialmente). \\ \hline
\end{tabularx}

\end{document}
